\section{La banda \'util real de los datos de campo}
\label{sec:banda-util}

\emph{An\'alisis nuevo del 12 de agosto de 2026.} Es la pregunta que el proyecto
entero intenta responder y que, con datos propios, nunca se hab\'ia contestado
directamente: \textbf{?`de qu\'e frecuencias dispone realmente la se\~nal
registrada en campo?} Toda la cadena anal\'ogica ---la compensaci\'on, la
ganancia agresiva, el DA-AFE que la ganancia obliga--- existe para recuperar el
extremo bajo. Corresponde ver cu\'anto se recuper\'o.

\subsection{M\'etodo}

Sobre 265 golpes de la campa\~na Canchita en 15 distancias, y para cada golpe:

\begin{itemize}[itemsep=0.15em]
    \item \textbf{Ventana de se\~nal}: \SI{1}{\second} desde el instante de
    disparo del martillo.
    \item \textbf{Ventana de ruido}: \SI{1}{\second} \emph{anterior} al disparo.
    Los registros arrancan aproximadamente \SI{1.24}{\second} antes del
    \emph{trigger}, de modo que esa ventana existe en todas las capturas y
    contiene el ruido ambiente y de instrumento en las condiciones exactas de la
    medici\'on.
    \item Densidad espectral de potencia de ambas por ventana de Hann, mediana
    por distancia, y
    $\mathrm{SNR}(f)=10\log_{10}\!\left(P_{\text{se\~nal}}/P_{\text{ruido}}\right)$.
\end{itemize}

Los datos se toman sin filtrar (el estado guardado del servidor es
\emph{sin filtrar}, orden 4, cortes en cero), de modo que lo que se ve es lo que
entreg\'o la cadena.

\begin{figure}[ht]
\centering
\includegraphics[width=\textwidth]{contenido_baja_frecuencia.png}
\caption{Arriba: densidad espectral de la se\~nal por distancia. Abajo: SNR
contra el ruido previo al disparo. La banda sombreada es el objetivo del proyecto
(\SIrange{1}{2.5}{\hertz}, correspondiente a \SI{50}{\meter} de profundidad); la
vertical punteada es la frecuencia natural del SM-24.}
\label{fig:banda-util}
\end{figure}

\subsection{El resultado}

\begin{center}\footnotesize
\begin{tabular}{@{}lrrl@{}}
\toprule
Banda & SNR mediana & Grupos con SNR${}>\SI{10}{\decibel}$ & Veredicto \\
\midrule
\textbf{\SIrange{1}{2.5}{\hertz}} & \textbf{\SI{0.4}{\decibel}} & \textbf{0 de 16} & \textbf{La banda objetivo: sin se\~nal} \\
\SIrange{2}{5}{\hertz}   & \SI{1.2}{\decibel}  & 0 de 16 & Sin se\~nal utilizable \\
\SIrange{5}{10}{\hertz}  & \SI{1.3}{\decibel}  & 0 de 16 & Sin se\~nal utilizable \\
\SIrange{10}{20}{\hertz} & \SI{15.2}{\decibel} & 13 de 16 & Utilizable \\
\SIrange{20}{50}{\hertz} & \SI{25.6}{\decibel} & 14 de 16 & Utilizable \\
\bottomrule
\end{tabular}
\end{center}

\emph{Los grupos son combinaciones (distancia, $F_s$) con al menos tres golpes, de modo
que no se mezclan tasas de muestreo distintas dentro de un mismo promedio.}

\dato{\textbf{La banda \'util efectiva de las campa\~nas realizadas es
aproximadamente \SIrange{10}{50}{\hertz}}. Por debajo de \SI{10}{\hertz} ---es
decir, por debajo de la frecuencia natural del transductor--- la relaci\'on
se\~nal--ruido es de \emph{cero decibel}: la se\~nal y el ruido previo al golpe
tienen la misma potencia.}

\dato{\textbf{El diagn\'ostico que lo confirma}: en el panel superior, la
densidad espectral por debajo de \SI{7}{\hertz} es \emph{la misma para todas las
distancias}, plana en torno a \SIrange{-60}{-65}{\decibel}. Una onda que se
propaga tiene que atenuarse con la distancia; algo que no cambia entre
\SI{10}{} y \SI{50}{\meter} \textbf{no est\'a viajando por el suelo}. Es un piso
com\'un ---instrumental o ambiental--- y no el campo de ondas. Por encima de
\SI{10}{\hertz}, en cambio, las curvas se separan ordenadamente por distancia,
que es exactamente lo que debe hacer una se\~nal s\'ismica.}

\subsection{Controles de robustez}

El resultado es lo bastante consecuente como para exigir verificaci\'on. Se
hicieron tres controles, y los tres lo sostienen.

\paragraph{Control 1: ?`lo borr\'o el procesamiento?} No. Se repiti\'o el
c\'alculo sobre los \textbf{CSV crudos} de \texttt{data/raw/Canchita}, que nunca
pasaron por el \emph{pipeline}. Tres capturas independientes a
\SI{2929}{\hertz} dan densidades espectrales de \SIrange{1.5}{6.8}{\decibel} en
\SIrange{1}{5}{\hertz} contra \SIrange{22.7}{23.5}{\decibel} en
\SIrange{10}{20}{\hertz}: los mismos \textasciitilde{}\SI{17}{\decibel} de
diferencia. Las bajas frecuencias no se perdieron al procesar; no estaban.

\paragraph{Control 2: ?`es un artefacto de mezclar las dos configuraciones?} No.
Separando por configuraci\'on y por frecuencia de muestreo:

\begin{center}\footnotesize
\begin{tabular}{@{}llrrrrr@{}}
\toprule
Config. & $F_s$ & $n$ & \SIrange{1}{5}{\hertz} & \SIrange{5}{10}{\hertz} & \SIrange{10}{20}{\hertz} & \SIrange{20}{50}{\hertz} \\
\midrule
A & \SI{1020}{\hertz} & 12  & \SI{3.0}{\decibel} & \SI{0.8}{\decibel} & \SI{5.1}{\decibel}  & \SI{11.5}{\decibel} \\
A & \SI{2929}{\hertz} & 20  & \SI{2.0}{\decibel} & \SI{1.3}{\decibel} & \SI{13.8}{\decibel} & \SI{26.1}{\decibel} \\
B & \SI{1020}{\hertz} & 238 & \SI{0.4}{\decibel} & \SI{1.5}{\decibel} & \SI{21.6}{\decibel} & \SI{24.4}{\decibel} \\
\bottomrule
\end{tabular}
\end{center}

Las tres combinaciones coinciden: \SIrange{0.4}{3.0}{\decibel} por debajo de
\SI{10}{\hertz}, mucho m\'as arriba de \SI{10}{\hertz}.

\paragraph{Control 3: ?`la ventana de \SI{1}{\second} alcanza para medir
\SIrange{1}{2.5}{\hertz}?} Es la objeci\'on m\'as seria ---\SI{1}{\second} son
apenas una a dos oscilaciones a esas frecuencias--- y merece respuesta directa.
Repitiendo con ventanas de se\~nal m\'as largas:

\begin{center}\footnotesize
\begin{tabular}{@{}lrrrr@{}}
\toprule
Ventana de se\~nal & \SIrange{1}{2.5}{\hertz} & \SIrange{5}{10}{\hertz} & \SIrange{10}{20}{\hertz} & \SIrange{20}{50}{\hertz} \\
\midrule
\SI{1}{\second} & \SI{0.4}{\decibel}  & \SI{1.3}{\decibel}  & \SI{15.2}{\decibel} & \SI{25.6}{\decibel} \\
\SI{2}{\second} & \SI{-0.1}{\decibel} & \SI{0.3}{\decibel}  & \SI{15.1}{\decibel} & \SI{15.2}{\decibel} \\
\SI{4}{\second} & \SI{-0.1}{\decibel} & \SI{-0.1}{\decibel} & \SI{5.3}{\decibel}  & \SI{4.5}{\decibel} \\
\SI{6}{\second} & \SI{-0.6}{\decibel} & \SI{-0.5}{\decibel} & \SI{1.7}{\decibel}  & \SI{1.6}{\decibel} \\
\bottomrule
\end{tabular}
\end{center}

\seleccionado{\textbf{La conclusi\'on en baja frecuencia es independiente de la
ventana}: \SIrange{1}{2.5}{\hertz} se queda entre \SI{-0.6}{} y
\SI{+0.4}{\decibel} en los cuatro casos. Es la forma m\'as fuerte del hallazgo.

La ca\'ida de la SNR en alta frecuencia al alargar la ventana no es un problema
sino f\'isica esperable: el paquete s\'ismico dura una fracci\'on de segundo, de
modo que extender la ventana agrega sobre todo silencio posterior y diluye la
potencia de se\~nal. Confirma, de paso, que \textbf{la ventana de
\SI{1}{\second} es la elecci\'on correcta}, porque es la que se ajusta a la
duraci\'on del evento.}

\paragraph{Control 4: ?`es propio de un solo conjunto procesado?} No. Repetido
sobre \texttt{Canchita\_grupo1\_procesado}, un conjunto procesado independiente
(250 golpes, 14 grupos), da el mismo cuadro: \SI{1.0}{\decibel} en
\SIrange{1}{2.5}{\hertz} (0 de 14 grupos por encima de \SI{10}{\decibel}),
\SI{1.7}{\decibel} en \SIrange{5}{10}{\hertz}, \SI{17.2}{\decibel} en
\SIrange{10}{20}{\hertz} (11 de 14) y \SI{24.9}{\decibel} en
\SIrange{20}{50}{\hertz} (12 de 14).

\pendiente{\textbf{Un l\'imite que los datos imponen}: la ventana de ruido no
puede superar \textasciitilde{}\SI{1.2}{\second}, porque es todo lo que los
registros graban antes del disparo. No es posible, con este conjunto, comparar
ventanas iguales m\'as largas. Para una medici\'on limpia del piso de ruido a
muy baja frecuencia hace falta una \textbf{captura larga sin golpe}, que el
sistema ya sabe hacer (la captura continua de 10 minutos existe) pero que nunca
se us\'o con este fin.}

\subsection{Qu\'e significa y qu\'e no significa}

\seleccionado{\textbf{Lo que no significa.} No dice que la compensaci\'on sea
in\'util ni que el enfoque sea equivocado. Dice que \emph{en las campa\~nas
efectivamente realizadas} el extremo bajo no lleg\'o a estar por encima del
ruido. Y eso es coherente con lo que ya se sab\'ia del banco: la muesca de
cancelaci\'on medida qued\'o \textbf{\SI{26.33}{\decibel} menos profunda} que el
modelo (\S\ref{sec:atlas}), es decir que la compensaci\'on realizada \emph{no
alcanz\'o} el $\zeta_1$ nominal. Las dos mediciones ---banco y campo--- apuntan
en el mismo sentido y se corroboran mutuamente.}

\pendiente{\textbf{Lo que s\'i significa, y es incómodo pero \'util.} El rango
\SIrange{10}{50}{\hertz} que los datos s\'i tienen es aproximadamente
\textbf{el mismo rango en el que Park \emph{et al.} operaban sin necesidad de
compensar}. Dicho de frente: hasta ahora, las campa\~nas de este proyecto
midieron en el r\'egimen del MASW somero. La contribuci\'on que la
compensaci\'on busca ---el r\'egimen por debajo de $f_n$--- \textbf{todav\'ia no
est\'a demostrada en campo}, aunque el mecanismo est\'e implementado y medido en
banco.}

\seleccionado{\textbf{Y corrobora, desde una segunda direcci\'on independiente,
el problema de los \emph{picks} de baja frecuencia.} Los tres puntos de la curva
de dispersi\'on con $\lambda>\SI{40}{\meter}$ (\S\ref{sec:contraste-sev}) est\'an
en \SIrange{5}{7}{\hertz}, es decir \textbf{en una banda con
\SIrange{0}{2}{\decibel} de SNR}. No es s\'olo que la apertura no los resuelva:
es que probablemente \emph{no hay se\~nal ah\'i}. Dos an\'alisis independientes
---geometr\'ia del arreglo y relaci\'on se\~nal--ruido--- se\~nalan los mismos
puntos. Deben excluirse de la inversi\'on.}

\subsection{Consecuencia para el orden de prioridades}

Este resultado reordena la lista de pendientes del proyecto. Hasta ahora el
cuello de botella declarado era la \textbf{apertura del arreglo}
(\SI{40}{\meter} contra el criterio $\lambda_{\max}>\SI{60}{\meter}$ de Comina).
Sigue si\'endolo, pero \textbf{no es el primero}:

\begin{enumerate}[itemsep=0.2em]
    \item \textbf{Primero, que exista se\~nal por debajo de \SI{10}{\hertz}.}
    Ampliar la apertura para resolver longitudes de onda largas no sirve si en
    esas frecuencias no hay nada por encima del ruido. Las v\'ias concretas, en
    orden de costo creciente:
    \begin{itemize}[itemsep=0.1em]
        \item \textbf{Punta m\'as blanda y masa de extensi\'on} en el martillo
        (\S\ref{sec:fuente-sismica}): mueve energ\'ia de la fuente hacia abajo,
        y no cuesta nada.
        \item \textbf{Terminar la calibraci\'on de la cancelaci\'on}: recuperar
        los \SI{26.33}{\decibel} de profundidad de muesca que faltan es,
        literalmente, el objetivo de dise\~no incumplido.
        \item \textbf{Apilar muchos m\'as golpes.} El apilado mejora la SNR con
        $\sqrt{N}$: pasar de \textasciitilde{}20 golpes por posici\'on a 200
        dar\'ia unos \SI{10}{\decibel}, que es la diferencia entre
        \SI{0.5}{\decibel} y algo utilizable. Es la v\'ia m\'as barata de todas
        y la m\'as compatible con el m\'etodo de apertura sint\'etica ya adoptado.
        \item \textbf{Caracterizar el piso de ruido} para saber si es del
        instrumento o del sitio ---dos problemas completamente distintos, y hoy
        no se sabe cu\'al es. Basta una captura larga con el ge\'ofono
        desacoplado del suelo, que se puede hacer en el laboratorio.
    \end{itemize}
    \item \textbf{Despu\'es, la apertura}, para que las longitudes de onda que
    entonces s\'i existan puedan resolverse.
    \item \textbf{Por \'ultimo, el contraste externo} contra el SEV~01 y, si se
    ejecuta, contra la perforaci\'on de \SI{60}{\meter}.
\end{enumerate}

\pendiente{Una salvedad honesta sobre este an\'alisis: la ventana de ruido es
\emph{previa al golpe} y por lo tanto no incluye el ruido inducido por el propio
impacto (\emph{ground roll} tard\'io, reverberaci\'on, ruido de acoplamiento).
La SNR real dentro de la ventana de se\~nal podr\'ia ser menor a\'un. Tambi\'en
conviene rehacerlo separando las dos configuraciones de adquisici\'on
identificadas en \S\ref{sec:repetibilidad-fuente}, y comprobar que el resultado
no cambie.}